\subsection*{20. 2-Part Tariffs, Minimum Commitments, and Hybrid Models}

An access fee plus usage price can split the value of availability and consumption. A minimum contractual commitment can finance dedicated capacity; additional usage can follow a transparent schedule. The architecture is attractive when the seller bears fixed readiness costs but the buyer's volume varies.

Yet complexity creates buyer risk. A customer who cannot explain a contract to finance may decline it regardless of theoretical efficiency. Every additional fee should have a genuine economic or operational function. Founders should simplify until the customer can predict the bill and the firm can reasonably forecast costs.

\subsection*{21. Freemium and Free Trials Are Financing Decisions}

Free access can reduce trial friction, support network effects, or let customers experience a product before committing. But ``free'' customers generate infrastructure, support, abuse, and acquisition costs, and may rarely convert. If the company serves a large free audience hoping for future paid upgrades, that future value must be measured, not assumed.

A freemium offering should have a coherent free use case and a credible reason to upgrade when the job becomes more demanding. A free trial should reveal the specific value that a buyer needs to assess, not be an arbitrary countdown. In education, free teaching can be genuine public service and a demonstration of pedagogy, but it cannot by itself prove that buyers will fund a high-touch program.

\subsection*{22. Licensing, Royalties, and Intellectual Property}

A company can capture value by licensing software, curriculum, data rights, designs, or technical methods. The economic benefit is reuse across customers without reproducing all development effort. The risks include legal scope, enforcement, versioning, quality control, customer dependence, and whether the intellectual property provides real value independent of services.

A college curriculum license may sound scalable until every client requires custom instructional delivery, integration, training, and assessment. Measure the true reusable share. Royalty structures can align payments with activity but require reliable measurement and audits. The licensing price should reflect what the customer can deploy and gain, not simply what the author spent creating the material.

\subsection*{23. Outcome-Based and Performance-Linked Fees}

A performance fee attempts to align seller payment with buyer value. It can be appropriate when the outcome is observable, attributable, timely, auditable, and partly under the provider's control. In practice, outcomes may depend on the customer's staffing, baseline trends, external shocks, and reporting choices. Contracts then produce disputes over what counts as incremental performance.

A safer architecture often combines a base fee covering legitimate delivery with a bounded bonus for a clearly defined metric, verified under a pre-agreed method. Do not sell guaranteed employment, arbitrary salary increases, or unmeasurable enterprise transformation. Where attribution is weak, milestone-based fees for controllable deliverables may be more honest than ``pay only for results.''
